This point-to-point architecture choice ensures the resilience and privacy of 
telecommunications. However, authenticity is more difficult to preserve.

\subsection{Privacy}

The term \enquote{privacy} is semantically rich. It ranges from an individual's 
physical condition to their most fundamental rights \cite{rfc6973}. P2P 
connections avoid some threats. For instance, surveillance is made more tricky; 
uniquely if traffic is encrypted. Since data is stored locally, it relies on 
hardware security measures and only compromises one user in the event of a 
leak. However, it discloses the IP address, this gives an approximate 
geolocation of the peer. To maximize benefits, you SHOULD use onions, but it 
remains complicated.

\subsection{Authenticity}

Authenticity allows the author of a message to be proven. In P2P, authenticity 
is hard to prove. Turms allows individuals to authenticate themselves through 
an Autha instance, or to remain as \enquote{guests.}

\begin{enumerate}
    \item \textbf{Autha instance with discovery server}: discovery server 
validates the decentralized identity, then generates a JWT. During signaling, 
the discovery server acts as a trusted third party. User acquires the ID 
\enquote{user\_id@discovery\_server.tld.} Once the initial connection has been 
established, messages are authenticated by their signature.
    \item \textbf{Autha instance without discovery server}: user generates a 
token that will be sent to the peer, who will verify it using the public key 
provided by the user's Autha instance.
    \item \textbf{As guest}: untrustworthy and inauthentic by nature. Trust is 
acquired through the SDP transmission channel.
\end{enumerate}

\subsection{Integrity}

Integrity ensures that messages has not been altered during transit. This 
component goes hand in hand with data encryption.

\subsection{Resilience}

P2P simplifies resilience. By design, there is no longer a single point of 
failure, except for the network itself.